% ch1_overview.tex — 1장 개요
\section{개요}
\label{sec:overview}

\noindent{\sffamily\small\color{muted}\textbf{CHAPTER 1}}

\begin{quoteblock}
서울 25개 자치구의 디지털 격차를 측정 $\rightarrow$ 설명 $\rightarrow$ 복원의 3단계 분석 체계로 검토한다. 비판적 실재론의 층화된 존재론을 메타이론으로 두어, 측정값\textperiodcentered{}계수\textperiodcentered{}텍스트가 동일한 디지털 연결성 현상의 다른 단면임을 명시한다.
\end{quoteblock}

\figfull{image1.png}{UMC 6개 차원 측정 체계}{fig:0-1}
\figfull{image2.png}{3단계 분석 파이프라인}{fig:0-2}

\subsection{이론적 이해}
\label{subsec:1-1}

본 프로젝트는 서울의 디지털 연결성을 세 개의 분석 단계로 검토하고, 그 결과를 바탕으로 정책 제언을 도출한다. 디지털 연결성은 ITU가 제시한 UMC(Universal Meaningful Connectivity) 프레임에 기반한다. 프로젝트의 분석은 개별 결과물의 병치가 아니라 하나의 메타이론적 틀 안에서 해석된다. 메타이론은 현실의 본성에 관한 포괄적 관점으로서 연구와 실천을 철학적으로 뒷받침하며~\citep{AllanaClark2018}, 그 안에서 각 방법론의 존재론적 위치를 명시할 근거를 제공한다. 관측된 현상이 서로 다른 층위의 개념으로 일반화되는 과정을 명시적으로 드러냄으로써 프로젝트 자체의 타당성을 확보하고자 한다.

이러한 접근은 프로젝트 자체의 타당성 뿐만 아니라, 다른 맥락으로의 적용을 위한 명확한 근거를 제공한다는 이점이 있다. 이는 현재 정책 시행의 대안점을 제공할 수 있다. 증거기반 정책(evidence-based policy)이 보편화되면서 원래 맥락의 제도적\textperiodcentered{}문화적 조건을 충분히 고려하지 않고 분석 결과만을 차용해 적용하는 사례가 늘고 있다~\citep{DolowitzMarsh2000}. 한편, 증거기반 패러다임 자체에 대한 비판도 축적되고 있다. 좁은 의미의 증거기반 틀은 정책결정의 복잡하고 맥락 의존적이며 가치가 개입하는 속성을 포착하지 못한다는 지적이 대표적이다~\citep{GreenhalghRussell2009}. 이에 따라 증거기반 정책에서 증거참고 정책(evidence-informed policy)으로의 전환이 논의되고 있으며~\citep{Head2010}, 증거의 역할은 정책의 체계적 토대를 구성하기보다 정책을 안내하는 방향으로 재정립되고 있다.

본 프로젝트는 경험적 증거의 중요도를 낮추기보다, 이론적 틀에 기반해 다른 맥락에 적용되기 위한 조건과 핵심 구조를 정교하게 조명하고자 한다. 각 분석의 역할과 상호 연결, 그리고 정책 제언으로 이어지는 해석 구조를 명시적으로 제시하고, 해석을 위한 주요한 맥락 정보를 상세하게 밝히고자 한다. 이는, 개별 분석 방법과 구분되는 프로젝트 자체의 사례 연구적 성격을 고려한 것이다.

\begin{metabox}[title=메타이론 \textperiodcentered{} 층화된 존재론 --- 경험\textperiodcentered{}실재\textperiodcentered{}실제 영역]
관측된 디지털 격차(경험 영역)는 사회\textperiodcentered{}경제\textperiodcentered{}물리적 메커니즘(실재 영역)이 특정 조건에서 작동한 결과(실제 영역)이다. 본 분석은 세 영역을 명시적으로 분리하여 측정\textperiodcentered{}계수\textperiodcentered{}텍스트 산출물의 인식론적 위치를 다르게 둔다.
\end{metabox}

개별 분석 방법론은 세 가지로 구분된다. UMC의 이론적 프레임을 정책 시행 주체인 서울 25개 자치구에 적용하여 격차를 측정한다. 이어서 다층 모형을 이용해 개인의 디지털 연결 문제에 대한 개인 요인과 지역 요인의 기여를 분리한다. 마지막으로, 앞선 분석에서 일반화된 서울의 디지털 연결 문제에 대한 경험적 증거를 비정형 텍스트를 통해 관측된 개별 현상으로 복원한다. 이는, UMC의 개념적 틀에 맞게 개인의 경험을 추상화하여 재구성하는 절차를 포함한다. ITU가 제시한 UMC의 개념은 집계한 결과와 실제 경험적 근거를 비교한다. 마지막으로, 개념에서 실제 현상이 일어나는 인과적인 기재를 추론한다. 이러한 복합적인 층위 전환은 비판적 실재론이 제시하는 층화된 존재론, 즉 경험적 영역, 실재적 영역, 실제적 영역의 구분에 기반한다~\citep{Bhaskar1975}.

이 과정이 필요한 이유를 가정적 상황으로 설명한다. 예를 들어 관측된 경험적 증거를 통해, 급행 노선이 지나는 지하철 역의 역세권 거주지가 그렇지 않은 역의 거주지보다 주택 가격이 높다는 결과가 도출되었다고 하자. 결과의 타당성을 떠나 해당 결과 자체는 유의미한 지식으로 보기 어렵다. 이 결과는 서울에서만 적용될 수도 있다. 또한, 결과를 근거로 모든 역에 급행 노선이 지나도록 하는 것은 정책적으로도 논리적으로도 실현 가능하지 않다. 세계에 관한 보편타당한 지식은 경험적 근거의 일반화를 넘은, 타당한 추론을 통해 밝혀진다~\citep{Bhaskar1975}. 앞선 상황에서, 가령 서울의 시간 가치가 높기 때문에 시간을 보다 줄일 수 있는 급행 노선이 주택 가격에도 반영된다는 기제를 밝히는 것이 그 예일 것이다. 다음 절에서는 개별 분석 방법론 자체가 갖는 다른 맥락에서의 적용 이점을 설명한다.

\subsection{분석 방법 이해}
\label{subsec:1-2}

본 절의 목적은 세 가지 분석 방법을 설명하고, 앞서 제시한 메타이론적 틀과는 별도로 각 방법론 자체가 다른 맥락에 적용될 수 있는 근거를 밝히는 데 있다. 우선, 데이터에 대해 회귀 기반 분석을 통해 경험적인 패턴을 식별한다. 개별 관측을 측정하고 추상화한 뒤 다시 개별 현상으로 복원하며, UMC 개념이 경험적으로도 나타나는지 검증하는 메타 구조의 타당성과 개별 방법론의 적용 타당성을 분리하여 논의한다.

먼저, 서울 자치구별 UMC를 측정하여 어디에서 어떤 차원의 격차가 존재하는지 보여준다. 측정 지표는 ITU(2023)가 제시한 6개 차원의 개념적 정의에 기반하되, 서울의 데이터 환경과 도시 특성에 맞게 조작화하였다. 정책을 실제로 수행하는 주체는 국가뿐 아니라 지방정부를 비롯한 다양한 단위에 걸쳐 있다. 이에 본 프로젝트는 측정 및 분석의 단위를 정책 수행 단위에 맞도록 서울 25개 자치구로 설정하였다. 개념적 정의와 대상지에 적합한 조작적 정의를 분리하여 적용하였으므로, 동일한 6개 차원 구조를 유지하면서 각 도시의 데이터 환경에 맞게 지표를 대체할 수 있다. 측정 지표 선정과 결과는 \ref{subsec:3-1}절에서 상술한다.

그 다음, 개인 특성과 자치구 수준의 UMC 인프라 특성을 함께 투입한 2수준 위계적 선형 모형(HLM)을 통해, 디지털 활용의 차이가 누구에게 어떤 조건에서 크게 나타나는지 검토한다. 이는 디지털 활용 격차의 인과 기제를 경험적으로 일반화하기 위함이다. 동일 행정 지역 안에서도 개인의 교육 수준, 연령, 가구 구성에 따라 디지털 활용이 크게 달라질 수 있기 때문이다. HLM은 개인이 지역에 내포된 위계적 자료 구조를 고려하는 표준적 방법이며, 투입 변수를 해당 맥락의 인구 특성과 인프라 조건에 맞게 교체하면 동일한 분석 틀을 다른 도시에도 적용할 수 있다. 변수 선정, 모형 설정, 교차수준 상호작용의 구체적 설계는 \ref{subsec:3-2}절에서 상술한다.

마지막으로, 비정형 텍스트를 통해 앞선 두 분석이 포착하지 못하는 개별 현상을 복원하고, UMC 개념에 적합한 상위 구조로 추상화하며, 추론된 인과 기제와 관측된 경험 세계 간의 정합성을 검증한다. 개별 사례를 지역 수준의 UMC 개념과 곧바로 동일시하지 않고 추상화 과정을 거치며, 앞선 관측 결과와 비교하여 심층적인 인과 기제를 밝힌다. 이 과정이 필요한 이유를 가정적 상황으로 설명한다.

텍스트 분석에는 두 가지 요구가 있다. 첫째, 개별 텍스트 사례가 디지털 연결에 관련된 것인지 판단하고, UMC의 6개 차원이라는 사전 정의된 이론 틀에 따라 분류할 수 있어야 한다. 둘째, 분석 방법이 특정 언어에 강하게 묶이지 않아야 한다. UMC의 프레임워크가 포괄적이기에 텍스트 분석 역시 언어 비종속성이 필요하다. 이 두 조건을 충족하기 위해 대규모 언어 모델(LLM)을 이론 주도적 텍스트 분류 도구로 사용한다. LLM은 형태소가 아닌 문맥 수준에서 의미를 처리하므로 한국어 생활 텍스트의 비표준적 표현에 강건하며, 프롬프트에 UMC 차원을 직접 제시하여 분류 기준을 연구자가 통제할 수 있다. 동일한 프롬프트 구조를 다른 언어의 텍스트에 적용할 수 있다는 점에서 방법론적 이전 가능성이 확보된다.

구체적으로, \ref{subsec:3-1}절의 UMC 실측값을 사전 분포로, 텍스트 분류 결과를 우도로 결합하여 지역별 UMC를 베이지안으로 추정한다. 이어서, \ref{subsec:3-2}절의 HLM 결과가 허용하는 범위 내에서 구조와 경험 사이의 메커니즘을 가시화한다. 마지막으로, 개별 게시글과 지역의 UMC 개념 사이의 인과적 기제를 추론한다. 이때 귀추적, 전향적, 경로적 추론의 세 가지 독립된 추론층을 설계하여, 각 추론 에이전트가 서로의 출력을 참조하지 않은 상태에서 가설을 생성한 뒤 상위 판정 에이전트가 이를 종합한다. 추론층을 분리함으로써 단일 추론 경로의 편향을 억제하고, 가설 간 수렴과 발산을 관찰할 수 있다. 프롬프트 설계, 오분류 억제 전략, 베이지안 업데이트 및 추론 절차의 구체적 내용은 \ref{subsec:3-3}절에서 상술한다.

정리하면, 본 프로젝트의 구조 안에서 2장은 3장의 결과를 읽기 위한 해석적 맥락을 제공하고, \ref{subsec:3-3}절은 \ref{subsec:3-1}절과 \ref{subsec:3-2}절의 결과를 생활세계의 경험과 다시 연결하는 매개 단계로 기능한다. 세 분석은 각기 다른 질문에 답하되, 앞선 분석의 산출물이 다음 분석의 입력이 되는 연쇄적 구조를 이룬다(그림~\ref{fig:project-overview} 참조). 각 방법론은 개념 구조와 조작적 정의를 분리하고, 교체 가능한 요소와 고정되는 요소를 명시함으로써 서울 이외의 맥락에서도 동일한 분석 틀을 적용할 수 있는 조건을 갖춘다(그림~\ref{fig:project-overview} 참조).

\figfull{image3.png}{프로젝트 개괄}{fig:project-overview}
